% !TeX root = ../main.tex
\section{Подгруппы и смежные классы}

\task{56.32, г}{Все подгруппы $A_4$}
Группа $A_4$ состоит из $e$, восьми тройных циклов и трёх
двойных транспозиций. Поэтому возможные порядки её элементов --- $1,2,3$.
Подгруппы имеют следующий вид:
\[
\begin{gathered}
\{e\},\qquad A_4,\qquad
V_4=\{e,(12)(34),(13)(24),(14)(23)\},\\
\langle(12)(34)\rangle,\quad
\langle(13)(24)\rangle,\quad
\langle(14)(23)\rangle,\\
\langle(123)\rangle,\quad\langle(124)\rangle,\quad
\langle(134)\rangle,\quad\langle(234)\rangle.
\end{gathered}
\]
Всего $10$ подгрупп. Обоснуем полноту списка.
По Лагранжу порядок подгруппы делит $12$, то есть равен
$1,2,3,4,6$ или $12$. Подгруппа порядка $2$ или $3$ циклическая:
любой её неединичный элемент имеет соответствующий простой порядок.
Тройные циклы объединяются в четыре пары взаимно обратных,
что даёт четыре подгруппы порядка $3$.
В подгруппе порядка $4$ не может быть элемента порядка $3$;
все её три неединичных элемента должны быть двойными транспозициями,
поэтому она равна $V_4$.

Предположим, что есть $H$ порядка $6$. Её индекс равен $2$,
поэтому $H\trianglelefteq A_4$ (доказательство этого факта приведено
в задаче 58.7). Факторгруппа $A_4/H$ имеет порядок $2$.
Образ любого тройного цикла в ней одновременно имеет порядок,
делящий $3$ и $2$, поэтому равен $e$.
Значит $H$ содержит все восемь тройных циклов и единицу,
что невозможно при $|H|=6$.
\editnote{Отсутствие элемента порядка $6$ само по себе не исключает
 подгруппу порядка $6$: например, $S_3$ имеет порядок $6$, но не имеет
 элемента порядка $6$. В списке исходника $(14)(24)$ заменено на $(14)(23)$.}

\task{56.37}{Описание смежных классов}
\textbf{в)} В аддитивной группе $\R$ по подгруппе $\Z$ классы равны
\[x+\Z=\{x+k:k\in\Z\},\qquad 0\le x<1.\]
Каждый вещественный элемент имеет единственную дробную часть
в $[0,1)$; два элемента принадлежат одному классу ровно тогда,
когда их разность целая.

\textbf{д)} В $\C^*$ по подгруппе $U$ имеем
\[zU=|z|U=\{w\in\C^*:|w|=|z|\}.\]
Действительно, умножение на элемент модуля $1$ не меняет модуль.
Обратно, если $|w|=|z|$, то $w/z\in U$ и $w\in zU$.
Классы параметризуются числами $r>0$.

\textbf{з)} Пусть $H=\Stab(n)=\{\sigma\in S_n:\sigma(n)=n\}$.
Для левого класса
\[gH=\{\sigma\in S_n:\sigma(n)=g(n)\}.\]
Одно включение следует из $(gh)(n)=g(n)$.
Для обратного, если $\sigma(n)=g(n)$, то
$(g^{-1}\sigma)(n)=n$, откуда $g^{-1}\sigma\in H$.
Получаем ровно $n$ классов, по одному на каждый образ точки $n$.
Если нужны правые классы, то
\[Hg=\{\sigma\in S_n:\sigma^{-1}(n)=g^{-1}(n)\}.\]

\task{56.38}{Классы по $\SL_n(\C)$}
\utv{}{Для $g\in\GL_n(\C)$ и $H=\SL_n(\C)$
 \[gH=\{a\in\GL_n(\C):\det a=\det g\}.\]
}{
 Если $a=gh$, $h\in H$, то $\det a=\det g\det h=\det g$.
 Обратно, при $\det a=\det g$ имеем
 $\det(g^{-1}a)=(\det g)^{-1}\det a=1$,
 то есть $g^{-1}a\in H$ и $a\in gH$.
}

\task{56.40}{Включение смежных классов}
\utv{}{Для подгрупп $H_1,H_2\le G$
 \[g_1H_1\subseteq g_2H_2
 \quad\Longleftrightarrow\quad
 H_1\subseteq H_2\ \text{и}\ g_2^{-1}g_1\in H_2.\]
}{
 Пусть выполнено включение. Так как $e\in H_1$, элемент $g_1$
 принадлежит $g_2H_2$, поэтому $u=g_2^{-1}g_1\in H_2$.
 Для любого $h\in H_1$ также $g_1h\in g_2H_2$, то есть
 $uh\in H_2$. Тогда $h=u^{-1}(uh)\in H_2$.

 Обратно, если $H_1\le H_2$ и $u=g_2^{-1}g_1\in H_2$,
 то для $h\in H_1$ имеем $g_1h=g_2uh\in g_2H_2$.
}

\task{56.41}{Пересечение смежных классов}
\utv{}{Если $g_1H_1\cap g_2H_2\ne\varnothing$ и
 $x\in g_1H_1\cap g_2H_2$, то
 \[g_1H_1\cap g_2H_2=x(H_1\cap H_2).\]
}{
 Из $x=g_ih_i$, $h_i\in H_i$, следует $xH_i=g_iH_i$.
 Теперь для любого $y\in G$ имеем
 \[
 y\in xH_1\cap xH_2
 \Longleftrightarrow x^{-1}y\in H_1\cap H_2
 \Longleftrightarrow y\in x(H_1\cap H_2).
 \]
 Это доказывает равенство обоих множеств.
}

\task{56.43}{Подмножество с условием $xy^{-1}z\in K$}
\utv{}{Если $K\ne\varnothing$ и $xy^{-1}z\in K$
 для всех $x,y,z\in K$, то $K$ --- правый смежный класс
 по некоторой подгруппе.}{
 Выберем $k\in K$ и положим $H=Kk^{-1}$.
 Тогда $e=kk^{-1}\in H$.
 Для $h_1=xk^{-1}$, $h_2=yk^{-1}$ имеем
 \[h_1h_2^{-1}=xy^{-1}=(xy^{-1}k)k^{-1}\in H.\]
 Последнее включение следует из условия при $z=k$.
 Непустое множество, замкнутое относительно $h_1h_2^{-1}$,
 является подгруппой: при $h_1=e$ получаем обратные,
 а затем $h_1h_2=h_1(h_2^{-1})^{-1}\in H$.
 Итак, $H\le G$, а $K=Hk$.
}

\task{56.44}{Мультипликативность индексов}
\utv{}{Если $H_1\le H_2\le G$,
 $[G:H_2]=m$ и $[H_2:H_1]=n$ конечны, то $[G:H_1]=mn$.}{
 Выберем представителей $a_1,\dots,a_m$ левых классов $G/H_2$
 и $b_1,\dots,b_n$ левых классов $H_2/H_1$.
 Каждый элемент $G$ имеет вид $a_i b_j h$, $h\in H_1$.
 Значит классы $a_i b_jH_1$ покрывают $G$.
 Если $a_i b_jH_1=a_{i'}b_{j'}H_1$, то элементы этих классов
 лежат одновременно в $a_iH_2$ и $a_{i'}H_2$, поэтому $i=i'$.
 Сокращая $a_i$, получаем $b_jH_1=b_{j'}H_1$, откуда $j=j'$.
 Имеем ровно $mn$ различных классов.
}
\editnote{Это доказательство работает и для бесконечной $G$.
 Отношения $|G|/|H_i|$ в этом случае использовать нельзя.}

\task{Т.2}{Подгруппы взаимно простых индексов}
\thm{}{Если конечные индексы $[G:H_1]=m$, $[G:H_2]=n$
 взаимно просты, то $G=H_1H_2$. Группа $G$ может быть бесконечной.}{
 Положим $K=H_1\cap H_2$. Отображение множеств классов
 \[G/K\longrightarrow G/H_1\times G/H_2,
 \qquad gK\longmapsto(gH_1,gH_2)
 \]
 инъективно: совпадение обеих координат означает
 $g_2^{-1}g_1\in H_1\cap H_2$.
 Следовательно, $q=[G:K]\le mn$, в частности индекс конечен.
 По задаче 56.44 числа $m,n$ делят $q$.
 Из $\gcd(m,n)=1$ следует $mn\mid q$, поэтому $q=mn$.

 Снова по мультипликативности
 $[H_1:K]=q/m=n$.
 Отображение $H_1/K\to G/H_2$, $hK\mapsto hH_2$,
 инъективно по тому же критерию совпадения классов.
 Обе стороны имеют $n$ элементов, значит оно сюръективно.
 Для любого $g\in G$ существует $h_1\in H_1$ с
 $gH_2=h_1H_2$; отсюда $g=h_1h_2$ для некоторого $h_2\in H_2$.
}
\editnote{Произведение $H_1H_2$ до доказательства не обязано быть
 подгруппой, поэтому выражение $[G:H_1H_2]$ использовать заранее нельзя.
 Из равенства $mk=nl$ также нельзя сразу заключать $mk=mn$:
 необходима верхняя оценка $[G:K]\le mn$.}
